\subsubsection{Comparaison des décompositions OpenMP à grande largeur}
Les anti-diagonales d’une matrice de \(m\times n\) cellules comptent \(m+n-1\) niveaux intérieurs de largeur variable. NvPD traite \(m\) lignes de \(n\) colonnes, dont les calculs peuvent avancer en parallèle. Les deux méthodes visitent les mêmes cellules, avec des accès mémoire et un nombre de synchronisations différents.

Le panneau a) de la figure~\ref{fig:million} compare les deux implémentations initiales sur les quinze paires d’un million de bases. Chaque méthode utilise huit cœurs EPYC. Pour les cinq paires à 5\,\% d’éditions introduites, les anti-diagonales prennent en moyenne \(190{,}24\pm0{,}42\)\,s et NvPD \(310{,}26\pm1{,}13\)\,s.

Le panneau b) reprend les mêmes cinq paires avec le noyau NvPD qui sélectionne seize prédécesseurs par permutation AVX-512. Un processus par paire dispose de huit cœurs physiques. Les cinq processus s’exécutent simultanément sur un nœud. Après une chauffe par méthode, trois blocs alternent l’ordre des deux solveurs. La lecture du FASTA précède le chronométrage, qui inclut la construction de l’index NvPD. Pour chaque paire \(j\), le gain moyen vaut
\[
  \bar S_j=\frac{1}{3}\sum_{b=1}^{3}\frac{T_{\mathrm{diag},bj}}{T_{\mathrm{NvPD},bj}} .
\]
\chart{cpu/openmp_un_million.png}{Le panneau a) donne les durées des implémentations initiales. Dans le panneau b), chaque point résume trois gains mesurés sur une paire. Le trait épais indique l’écart-type et le trait fin les extrêmes. Le losange donne la moyenne sur cinq paires, avec son écart-type et son IC à 95\,\%.}{fig:million}{bench/cpu/openmp/million/paired/trial.cpp}
Le temps moyen sur les cinq paires vaut \(196{,}864\pm0{,}123\)\,s pour les anti-diagonales et \(14{,}790\pm0{,}471\)\,s pour NvPD vectorisé. Les écarts-types décrivent la dispersion entre les cinq moyennes par paire. La figure~\ref{fig:maqao-permutation} montre comment la permutation réunit les lectures des prédécesseurs.

\synthese{À un million de bases et 5\,\% d’éditions introduites, NvPD vectorisé divise le temps des anti-diagonales par 13,322 sur huit cœurs EPYC. L’IC à 95\,\% calculé entre les cinq paires va de 12,787 à 13,857.}
